L'aire et le hamiltonien modulaire forment une paire conceptuellement analogue au produit et au quotient dans le vide.\par

L'opérateur d'aire détermine l'occupation totale. Le hamiltonien modulaire conserve la provenance sectorielle de cette occupation.\par

Au bord apparaissent deux familles de canaux : \(90\) et \(120\). Chacune possède ses opérateurs d'occupation. L'opérateur d'aire est construit à partir de la somme :\par

\[
\frac{\widehat{\mathcal A}_{\rm phys}}{\ell_P^2}
=
\gamma_{\rm BI}a_0
(N_{90}+N_{120}).
\]

L'aire compte le nombre total d'excitations. Qu'une excitation provienne du canal \(90\) ou du canal \(120\), l'aire élémentaire est la même. Du point de vue géométrique, les deux contribuent au même total.\par

Le hamiltonien holographique modulaire introduit une pondération sectorielle supplémentaire :\par

\[
K_{\rm HHM}-cI
=
\Delta_4^{\rm mod}
(90N_{90}+120N_{120}).
\]

Il compte les occupations et pondère en outre chaque excitation par le canal dont elle provient.\par

C'est pourquoi l'aire et le hamiltonien commutent : tous deux lisent les mêmes opérateurs d'occupation et contiennent des informations complémentaires.\par

Le spectre de l'opérateur d'aire est invariant par rapport à la provenance sectorielle, tandis que le hamiltonien modulaire la distingue.\par

L'opérateur d'aire détermine le nombre total d'excitations. La valeur modulaire pondérée détermine leur distribution entre les deux généalogies. La conservation conjointe des deux observables reconstruit exactement la décomposition sectorielle.\par

C'est une réalisation particulièrement nette de l'idée holographique. Le bord publie un ensemble minimal d'observables qui permet de récupérer l'information pertinente de l'intérieur.\par

L'holographie consiste ici à identifier quelle paire d'observables de frontière reconstruit exactement quelles données internes.\par

L'opérateur d'aire publie la grandeur géométrique et le hamiltonien modulaire publie la généalogie sectorielle. La paire d'observables reconstruit l'information complète des secteurs \(90/120\).\par

L'état réduit du bord transforme en outre ce hamiltonien en un véritable générateur modulaire :\par

\[
K_{\rm HHM}=-\log\rho_A.
\]

La purification en double champ thermique montre que l'entropie de l'état réduit est exactement l'entropie d'intrication au sein du système dupliqué déclaré. Il existe un état pur étendu dont la trace partielle produit l'état de bord et réalise explicitement l'identité thermique.\par

Apparaît alors le déficit informationnel orienté :\par

\[
I_{\rm or}
=
36\log3-S(\rho_A).
\]

Les \(36\) trits représenteraient la capacité maximale si tous les états étaient uniformes. La différence mesure la quantité d'information orientée conservée par la distribution réelle des poids \(90/120\).\par

L'entropie totale quantifie l'incertitude. Le déficit par rapport au maximum quantifie la structure orientée résiduelle.\par

La première loi modulaire traduit cette structure en géométrie :\par

\[
\delta S
=
\beta_{90}\,
\delta\langle\mathcal A_{90}\rangle
+
\beta_{120}\,
\delta\langle\mathcal A_{120}\rangle,
\]

con\par

\[
\frac{\beta_{120}}{\beta_{90}}
=
\frac43.
\]

Le changement d'entropie dépend conjointement du changement d'aire totale et du secteur de provenance. Deux variations ayant le même incrément géométrique peuvent avoir des réponses informationnelles différentes si elles proviennent de canaux différents.\par

Cette dépendance établit la conservation géométrique de la mémoire.\par

Pour les variations finies, l'entropie relative apparaît comme un terme positif de correction. L'extension exacte de la loi incorpore la différence entre la variation exacte et son approximation linéaire au moyen d'une quantité positive.\par

Cette correction positive détermine, hors du régime infinitésimal, le coût non linéaire de l'éloignement par rapport à l'état de référence et quantifie la contribution de mémoire généalogique dans la variation finie.\par

Enfin, la relation\par

\[
E_P=k_B\Theta_{\rm clk}\log3
\]

et le quantum d'aire\par

\[
\delta\mathcal A
=
\gamma_{\rm BI}a_0\ell_P^2
\]

permettent de définir une tension de conversion entre information énergétique et géométrie :\par

\[
\mathcal T_{I/A}
=
\frac{E_P}{\delta\mathcal A}.
\]

On exprime le rapport exact entre le quantum énergétique informationnel et le quantum géométrique d'aire dans la normalisation déclarée.\par

C'est pourquoi le hamiltonien holographique modulaire est une pièce centrale au sein d'un atlas de lecteurs complémentaires. Sa fonction spécifique consiste à conserver la provenance que la géométrie totale intègre dans une somme.\par

L'opérateur d'aire détermine l'occupation totale.\par
Le hamiltonien modulaire détermine la provenance sectorielle.\par
L'entropie quantifie l'information publiée.\par
L'entropie relative quantifie le coût de la transition généalogique.\par
